% ============================================================
%  cap00_apresentacao.tex
%  Apresentação (seção não numerada — frontmatter)
% ============================================================

\chapter*{Apresentação}
\addcontentsline{toc}{chapter}{Apresentação}

Este guia integra a série \textit{Avaliação na Prática}, publicada pelo FGV CLEAR --- Centro de Aprendizado e Pesquisa em Avaliação e Resultados, vinculado à Escola de Economia de São Paulo da Fundação Getulio Vargas (FGV EESP). A série tem como objetivo oferecer materiais de referência práticos e acessíveis para gestores, pesquisadores e técnicos envolvidos na avaliação de políticas públicas no Brasil.

Os volumes anteriores da série trataram de métodos de avaliação --- diferenças em diferenças, regressão descontínua, \textit{matching}, poder estatístico e análise de custo-benefício. Este guia ocupa um lugar diferente: em vez de discutir \textit{como} avaliar, ocupa-se de discutir \textit{com quais dados} avaliar. É um atlas do ecossistema de fontes de dados disponíveis para a avaliação de políticas públicas no Brasil.

\medskip

O Brasil é, sob muitos aspectos, um caso privilegiado para a avaliação quantitativa de políticas públicas. O país dispõe de um sistema estatístico nacional amadurecido, liderado pelo IBGE, de registros administrativos abrangentes consolidados ao longo de décadas, e de uma cultura crescente de abertura de dados que colocou as informações governamentais ao alcance de qualquer pesquisador com acesso à internet. O Censo Demográfico, a PNAD Contínua, o Censo Escolar, o DATASUS, a RAIS, o CadÚnico --- cada uma dessas fontes representou, em seu tempo, um avanço significativo na capacidade de compreender e avaliar as políticas públicas brasileiras.

Ao mesmo tempo, a abundância de fontes traz seus próprios desafios. Cada base foi criada com uma finalidade específica, obedece a sua própria lógica de produção e tem limitações que só se revelam a quem a conhece bem. Usar o Censo Escolar para uma pergunta que exige a PNAD, ou confundir os dados do SIH com uma medida da saúde da população em geral, são erros que comprometem análises tecnicamente sofisticadas. Este guia pretende ser o mapa que evita esses desvios.

\medskip

O guia está organizado em treze capítulos. O primeiro discute por que a escolha das fontes de dados é uma decisão metodológica com consequências diretas sobre a validade das avaliações e apresenta os critérios para selecionar fontes adequadas. Os capítulos seguintes percorrem as principais áreas temáticas --- tipologia de fontes, pesquisas amostrais e censos, mercado de trabalho, saúde, educação, assistência social, segurança pública, infraestrutura e território, finanças públicas e fontes internacionais --- antes de dedicar um capítulo às ferramentas tecnológicas emergentes (APIs, \textit{web scraping} e inteligência artificial) e encerrar com orientações práticas sobre como escolher, acessar, linkar e documentar fontes de dados.

Os Anexos A e B oferecem ferramentas de referência rápida: uma tabela-catálogo consolidada de todas as fontes apresentadas e um glossário dos principais termos técnicos.

\medskip

Este guia foi elaborado pela equipe do FGV CLEAR. Qualquer erro ou omissão é de nossa responsabilidade.

\bigskip

\noindent\textit{São Paulo, 2024}

\noindent\textbf{FGV CLEAR}
